\documentclass[11pt,a4paper]{article}

\usepackage[utf8]{inputenc}
\usepackage[T1]{fontenc}
\usepackage[french]{babel}
\usepackage[margin=2.5cm]{geometry}
\usepackage{graphicx}
\usepackage{hyperref}
\usepackage{listings}
\usepackage{xcolor}
\usepackage{amsmath}
\usepackage{enumitem}

\lstset{
    basicstyle=\ttfamily\footnotesize,
    breaklines=true,
    frame=single,
    backgroundcolor=\color{gray!10},
    language=C++,
    showstringspaces=false
}

\title{Projet final -- CSC\_43043\_EP : Île tropicale\\
\large Modélisation géométrique et programmation graphique}
\author{% TODO: remplacer par les noms des membres de l'équipe
Nom Prénom 1 \and Nom Prénom 2}
\date{\today}

\begin{document}

\maketitle

\section{Introduction}

Ce projet implémente une scène 3D interactive représentant une île
tropicale, développée avec la bibliothèque CGP utilisée dans le cours
de programmation graphique (CSC\_43043\_EP). La scène comprend un
terrain généré procéduralement (avec une crevasse / faille centrale),
un océan animé avec reflets du ciel, un ciel sous forme de skybox, une
forêt de palmiers et d'arbres tropicaux affichés par \emph{instancing}
GPU, une population d'oiseaux animés survolant la forêt, ainsi qu'un
bateau, une épave et un coffre au trésor posés sur le fond marin.

Ce rapport présente les principales techniques mises en \oe{}uvre, en
les reliant aux notions vues en cours (rasterisation, modèle
d'illumination de Phong, texturage, \emph{shadow mapping},
\emph{environment mapping}, bruit procédural et \emph{instancing} GPU),
ainsi que les difficultés techniques rencontrées et les choix
d'optimisation effectués pour obtenir une scène jouable en temps réel.

\section{Vue d'ensemble de la scène}

La scène est construite autour des éléments géométriques suivants :

\begin{itemize}
    \item un \textbf{terrain} (\texttt{mesh\_primitive\_grid}, grille de
    $200 \times 200$ sommets sur une surface de $50\times 50$ unités),
    déformé procéduralement pour former une île avec une faille
    centrale (section~\ref{sec:crevasse}) ;
    \item un plan d'\textbf{eau} (même résolution de grille), animé en
    temps réel pour simuler des vagues (section~\ref{sec:vagues}) et
    réfléchissant le ciel (section~\ref{sec:reflets}) ;
    \item un \textbf{ciel} sous forme de sphère texturée englobant toute
    la scène, faisant office de \emph{skybox} (section~\ref{sec:skybox}) ;
    \item une \textbf{végétation} composée de palmiers et d'arbres
    tropicaux, dupliqués en grand nombre par \emph{instancing} GPU
    (section~\ref{sec:vegetation}) ;
    \item des \textbf{oiseaux} animés qui décrivent des trajectoires en
    orbite autour de la forêt (section~\ref{sec:animation}) ;
    \item un \textbf{bateau}, une \textbf{épave} et un \textbf{coffre au
    trésor}, posés sur le fond marin pour enrichir la scène ;
    \item un système d'\textbf{ombres portées} calculé par
    \emph{shadow mapping} (section~\ref{sec:ombres}).
\end{itemize}

Toutes les surfaces sont éclairées avec le modèle d'illumination de
Phong (composantes ambiante, diffuse et spéculaire), comme vu dans le
cours, à travers le couple de shaders
\texttt{shaders/mesh/mesh.vert.glsl} et
\texttt{shaders/mesh/mesh.frag.glsl}.

\section{Génération du terrain et la crevasse}
\label{sec:crevasse}

Le terrain part d'une grille plane (\texttt{mesh\_primitive\_grid})
dont chaque sommet est déplacé verticalement par la fonction
\texttt{deform\_terrain()} (\texttt{src/scene.cpp}). Cette fonction
combine plusieurs couches de \textbf{bruit de Perlin}
(\texttt{noise\_perlin}), une notion directement issue du cours sur le
bruit procédural, pour façonner le relief :

\begin{enumerate}
    \item \textbf{Masque de continent} : la position $(x,y)$ de chaque
    sommet est d'abord ``déformée'' (\emph{domain warping}) par du bruit
    de Perlin, puis on calcule un masque gaussien
    $\exp(-d^2/40)$ qui définit la silhouette générale de l'île.
    \item \textbf{La faille (crevasse)} : une seconde couche de bruit de
    Perlin distord une ligne suivant l'axe $y$
    (\texttt{fault\_distortion}). La distance d'un point à cette ligne
    est passée dans une fonction \texttt{smoothstep(0.8, 2.5, dist)} :
    en dessous de $0.8$, le point est considéré comme appartenant au
    \emph{canyon} (donc immergé), au-delà de $2.5$ il appartient
    pleinement au relief de l'île. Ce profil, multiplié au masque de
    continent, ``découpe'' littéralement l'île en deux par une gorge
    sinueuse : c'est la \textbf{crevasse} centrale de la scène.
    \item \textbf{Relief rocheux (fBm)} : une somme de 4 octaves de
    bruit de Perlin en valeur absolue (\emph{fractional Brownian
    motion}, $|noise|$ pour obtenir des arêtes anguleuses) module
    l'altitude des falaises et sommets.
    \item \textbf{Calcul final de l'altitude} : à l'intérieur de l'île
    (\texttt{island\_mask > 0.01}), l'altitude vaut
    $z = \text{niveau\_mer} + \text{island\_mask}\times h_{\max} -
    \text{mountain\_noise}\times \text{island\_mask}$, avec
    $h_{\max}=8$. À l'extérieur (dans le canyon ou l'océan), l'altitude
    redescend vers le fond marin ($z=-5$), légèrement bruité pour éviter
    un fond plat.
\end{enumerate}

Le résultat est une île massive coupée en son centre par une gorge
profonde aux parois quasi verticales, donnant un fort effet de vertige
lorsqu'on survole le bord de la faille -- d'où son nom de
\emph{crevasse}. Les normales du maillage sont ensuite recalculées
(\texttt{m.normal\_update()}) pour que l'éclairage de Phong réponde
correctement à ce relief escarpé.

\section{Animation des vagues}
\label{sec:vagues}

Le plan d'eau est représenté par un second \texttt{mesh\_primitive\_grid}
($200\times 200$, même emprise que le terrain), stocké côté CPU dans
\texttt{water\_mesh}. À chaque image, la fonction \texttt{idle\_frame()}
parcourt l'ensemble des sommets et met à jour leur composante $z$ :

\begin{lstlisting}
for (int k = 0; k < water_mesh.position.size(); ++k) {
    vec3& p = water_mesh.position[k];
    float wave1 = 0.20f * std::sin(0.6f * p.x + 1.2f * time);
    float wave2 = 0.15f * std::sin(0.8f * p.y + 1.7f * time + 0.5f);
    float noise = 0.15f * noise_perlin({ p.x*0.2f + time*0.1f,
                                          p.y*0.2f + time*0.1f });
    p.z = water_base_z + wave1 + wave2 + noise;
}
water_mesh.normal_update();
\end{lstlisting}

La houle est donc la superposition de \textbf{deux ondes
sinusoïdales} se propageant selon $x$ et $y$ avec des fréquences et
phases différentes (pour éviter un motif trop régulier), à laquelle on
ajoute une couche de \textbf{bruit de Perlin animé dans le temps}
(le décalage \texttt{time*0.1} fait ``dériver'' le bruit, donnant
l'impression d'un clapot irrégulier en plus de la houle principale).

Les normales sont ensuite recalculées à partir des nouvelles positions,
ce qui est essentiel : ce sont elles qui pilotent à la fois
l'éclairage de Phong de l'eau et le calcul du reflet du ciel
(section~\ref{sec:reflets}).

\paragraph{Mise à jour GPU sans fuite mémoire.} Plutôt que d'appeler
\texttt{water.initialize\_data\_on\_gpu(water\_mesh)} à chaque image
(ce qui réalloue un nouveau VBO sans libérer l'ancien et provoque une
fuite mémoire GPU et un avertissement \texttt{cgp}), les nouvelles
positions et normales sont envoyées directement dans les buffers déjà
alloués via :
\begin{lstlisting}
water.vbo_position.update(water_mesh.position);
water.vbo_normal.update(water_mesh.normal);
\end{lstlisting}
Cette mise à jour \emph{en place} correspond à la bonne pratique
indiquée par \texttt{cgp} pour l'animation de maillages et a permis de
corriger une importante chute de performance observée initialement.

\section{Ciel et \emph{skybox}}
\label{sec:skybox}

Le ciel est représenté par une grande sphère
(\texttt{mesh\_primitive\_sphere}) de rayon $1$, mise à l'échelle d'un
facteur $100$ (\texttt{sky.model.scaling = 100.f}) afin d'englober
toute la scène et de rester visuellement ``à l'infini'' quels que
soient les déplacements de la caméra. Cette sphère est texturée par
l'image \texttt{assets/sky.jpg} appliquée en projection
\textbf{équirectangulaire} (coordonnées UV de la sphère $\to$ image
panoramique du ciel).

Pour que la sphère apparaisse comme un arrière-plan uniforme et non
comme un objet éclairé, ses coefficients de Phong sont réglés à
\texttt{ambient = 1.0}, \texttt{diffuse = 0.0}, \texttt{specular = 0.0} :
seule la couleur de la texture est restituée, sans ombrage --
exactement le principe d'une \emph{skybox} vu en cours (un environnement
englobant non affecté par l'éclairage de la scène, qui donne l'illusion
d'un horizon lointain).

La même image est chargée une seconde fois dans \texttt{texture\_sky},
cette fois utilisée comme \textbf{texture d'environnement} pour le
calcul des reflets sur l'eau (section suivante).

\section{Reflets du ciel sur l'eau}
\label{sec:reflets}

Le calcul des reflets est effectué dans
\texttt{shaders/mesh/mesh.frag.glsl}, dans la branche activée par
\texttt{is\_water == 1}. Pour chaque fragment de la surface de l'eau :

\begin{enumerate}
    \item on calcule le \textbf{vecteur de vue} $V$ (de la surface vers
    la caméra) et le \textbf{vecteur réfléchi} $R = \text{reflect}(-V,
    N)$ par rapport à la normale $N$ du fragment -- exactement la
    fonction \texttt{reflect()} d'\emph{environment mapping} vue en
    cours ;
    \item ce vecteur $R$ est converti en coordonnées sphériques
    $(\varphi,\theta)$ puis en UV par projection équirectangulaire,
    pour échantillonner \texttt{sky\_texture} (= \texttt{texture\_sky})
    : on ``regarde'' ainsi le ciel dans la direction réfléchie ;
    \item un terme de \textbf{Fresnel} approché,
    $f = (1-\max(V\cdot N,0))^3$, augmente la part de réflexion lorsque
    le regard est rasant par rapport à la surface (effet observé sur
    une vraie étendue d'eau : on voit le fond de près, le ciel se
    reflète à l'horizon) ;
    \item la couleur finale de l'eau est un mélange
    (\texttt{mix}) entre la couleur diffuse/Phong de l'eau et la
    couleur du ciel échantillonnée, pondéré par le facteur de Fresnel et
    par un coefficient global \texttt{reflectivity} (réglé à $0.85$,
    proche d'un miroir).
\end{enumerate}

Comme les normales de l'eau varient à chaque image sous l'effet des
vagues (section~\ref{sec:vagues}), la direction $R$ varie elle aussi,
ce qui fait ``onduler'' le reflet du ciel à la surface de l'eau --
renforçant l'impression de mouvement.

Par ailleurs, lorsque la caméra passe sous le niveau de l'eau
(\texttt{camera\_position.z < water\_level\_b}), la couleur ambiante de
toutes les surfaces est mélangée avec une teinte bleue, simulant
l'atténuation de la lumière sous l'eau.

\paragraph{Remarque.} Une variable \texttt{reflec} a été conservée dans
le code comme uniforme générique préparatoire à des expérimentations
supplémentaires de réflexion ; elle n'est pas utilisée par le shader
par défaut (id=3), ce qui produit un avertissement \texttt{cgp}
inoffensif (\emph{``Try to send uniform variable [reflec] to a shader
that doesn't use it''}), sans impact sur le rendu.

\section{Ombres portées (\emph{shadow mapping})}
\label{sec:ombres}

Le rendu des ombres suit la technique classique du \emph{shadow
mapping} en deux passes, vue en cours :

\paragraph{1. Passe de profondeur (point de vue de la lumière).}
Une direction de lumière directionnelle fixe \texttt{light\_direction}
définit une caméra orthographique placée loin de la scène dans cette
direction. La matrice \texttt{light\_view\_projection} combine une
matrice de vue (translation vers la position de la lumière) et une
matrice de \textbf{projection orthographique} construite manuellement
(adaptée à une lumière directionnelle, sans perspective). La scène est
rendue depuis ce point de vue dans un \emph{framebuffer object} dédié
(\texttt{shadow\_fbo}, de résolution $1024\times 1024$, en mode
\texttt{depth} uniquement), à l'aide d'une paire de shaders minimalistes
(\texttt{shaders/shadow/shadow\_depth.vert.glsl} /
\texttt{.frag.glsl}) qui ne calculent que
\texttt{gl\_Position = light\_view\_projection * model * position} et
ne produisent aucune couleur -- seul le \emph{depth buffer} est
exploité.

\paragraph{2. Passe principale.} La texture de profondeur résultante est
liée à l'unité de texture 5 (\texttt{shadow\_map}) et la matrice
\texttt{light\_view\_projection} est envoyée comme uniforme. Dans
\texttt{mesh.frag.glsl}, pour chaque fragment :
\begin{enumerate}
    \item sa position monde est projetée dans l'espace de la lumière
    (\texttt{light\_space = light\_view\_projection * position}), puis
    ramenée en coordonnées de texture $[0,1]$ ;
    \item sa profondeur \texttt{current\_depth} est comparée à la
    profondeur stockée dans \texttt{shadow\_map} via un filtrage
    \textbf{PCF (Percentage-Closer Filtering) $3\times 3$} : on moyenne
    le résultat du test sur les 9 texels voisins, ce qui adoucit les
    bords des ombres et réduit l'aliasing ;
    \item un biais (\texttt{shadow\_bias = 0.007}) évite l'\emph{acné de
    surface} (auto-ombrage parasite dû à la précision limitée de la
    profondeur) ;
    \item si le fragment est occulté, sa couleur finale est multipliée
    par \texttt{shadow\_strength = 0.55} (l'ombre atténue la lumière
    sans la supprimer totalement).
\end{enumerate}

L'eau est explicitement exclue du calcul d'ombre
(\texttt{use\_shadow = 0} pendant son rendu) afin de ne pas perturber
le calcul des reflets.

\paragraph{Optimisation.} La résolution de la \emph{shadow map} a été
réduite de $2048$ à $1024$ pixels : sur le GPU intégré utilisé pour les
tests (Intel HD Graphics 630), cela divise par quatre le coût de la
passe de profondeur sans dégradation visible à l'échelle de la scène,
contribuant à maintenir un \emph{framerate} jouable malgré le grand
nombre d'objets dessinés deux fois (passe ombre + passe principale).

\section{Végétation : \emph{instancing} GPU et difficulté de
``tree2''}
\label{sec:vegetation}

\subsection{Principe de l'\emph{instancing}}

La forêt est composée de deux familles d'objets : des \textbf{palmiers}
(un seul maillage) et des \textbf{arbres tropicaux} (``tree2'',
décomposés en plusieurs parties, voir ci-dessous). Plutôt que
d'effectuer un appel de dessin (\emph{draw call}) par arbre -- ce qui
serait prohibitif pour des centaines d'arbres -- chaque type d'objet
est dessiné en une seule fois grâce au \textbf{\emph{GPU instancing}}
(\texttt{glDrawElementsInstanced} / \texttt{gl\_InstanceID}, vu en
cours).

Concrètement (\texttt{initialize\_instancing()} /
\texttt{draw\_part\_instanced()} dans \texttt{scene.cpp}) :
\begin{itemize}
    \item des VBOs supplémentaires sont attachés au VAO de chaque objet,
    aux \texttt{layout location} 4 à 8 du \emph{vertex shader}
    (\texttt{instance\_model\_col0..3} pour les 4 colonnes de la matrice
    modèle, \texttt{instance\_color} pour une teinte par instance) ;
    \item ces attributs ont un \emph{divisor} de 1
    (\texttt{glVertexAttribDivisor}) : ils changent une fois par
    \emph{instance} et non par sommet ;
    \item à chaque image, on remplit ces tableaux avec les matrices
    modèles (position, échelle, rotation aléatoire) de toutes les
    instances, on les envoie au GPU, puis un \textbf{unique} appel
    \texttt{glDrawElementsInstanced} dessine toutes les copies.
\end{itemize}

Dans \texttt{mesh.vert.glsl}, un booléen \texttt{use\_instancing}
sélectionne entre la matrice modèle uniforme classique et la matrice
modèle reconstruite à partir des 4 colonnes par instance.

\subsection{Placement procédural}

Les positions des palmiers (\texttt{tree\_positions}/\texttt{tree\_scales})
et des arbres ``tree2'' (\texttt{tree2\_positions}/\texttt{tree2\_scales})
sont tirées aléatoirement parmi les sommets du \texttt{terrain\_mesh}
déjà généré (section~\ref{sec:crevasse}), en ne conservant que les
points dont l'altitude et la pente correspondent à une zone de forêt
``plausible'' (au-dessus du niveau de la mer, hors falaises trop
abruptes). Chaque arbre reçoit en plus une échelle et une rotation
autour de l'axe vertical tirées aléatoirement, ce qui évite l'effet de
répétition visuelle typique de l'\emph{instancing} naïf.

\subsection{Difficulté rencontrée : ``tree2'' livré comme un seul
fichier OBJ multi-objets}

Le modèle d'arbre tropical (\texttt{trees9.obj}) a posé une difficulté
spécifique : contrairement au palmier, qui est un maillage unique,
\textbf{toutes les parties de l'arbre (tronc, écorce, branches,
feuillage...) sont regroupées dans un seul fichier OBJ}, sous la forme
de 8 sous-objets nommés à l'intérieur du fichier :
\texttt{Walnut\_L}, \texttt{Mossy\_Tr}, \texttt{Bark\_\_\_S},
\texttt{Bark\_\_\_0}, \texttt{Bottom\_T}, \texttt{Sonnerat},
\texttt{Bark\_\_\_1} et \texttt{Oak\_Leav}, totalisant à eux seuls
environ \textbf{740~000 sommets}.

Cela posait deux problèmes :
\begin{enumerate}
    \item \textbf{Texturage/matériaux différents par partie} : écorce,
    feuillage et branches mortes n'ont pas la même texture/couleur ; il
    fallait donc pouvoir dessiner chaque partie séparément avec son
    propre matériau, ce qui est impossible si l'OBJ est chargé comme un
    bloc unique ;
    \item \textbf{Instancing par partie} : pour appliquer
    l'\emph{instancing} décrit ci-dessus, chaque partie doit posséder
    son propre \texttt{mesh\_drawable} (son propre VAO/VBO), donc être
    \emph{séparée} avant l'envoi sur GPU.
\end{enumerate}

La solution retenue a été d'utiliser \texttt{mesh\_load\_file\_obj}
en conservant le découpage par \emph{nom d'objet} défini dans le
fichier OBJ (lignes \texttt{o <nom>}), produisant un
\texttt{std::vector<mesh\_drawable>}
\texttt{tree2\_parts} de 8 éléments. Chaque partie est ensuite
initialisée sur GPU et préparée pour l'\emph{instancing}
indépendamment :
\begin{lstlisting}
for (auto& part : tree2_parts) {
    initialize_instancing(part);
}
\end{lstlisting}
Au moment du dessin, les 8 parties sont dessinées avec les
\emph{mêmes} matrices de transformation par instance (un arbre =
8~appels \texttt{draw\_part\_instanced}, un par partie), de sorte que
chaque arbre reste visuellement cohérent (toutes ses parties sont
translatées/mises à l'échelle/tournées ensemble) tout en autorisant des
matériaux différents par partie.

\paragraph{Compromis de performance.} Avec $\sim$740~000 sommets par
arbre et 8~appels de dessin instanciés par passe (×2 passes pour
l'ombre), le coût en sommets traités croît très vite avec le nombre
d'instances. Sur le GPU intégré utilisé pour les tests, des essais à 8,
20 puis 150 instances ont montré une dégradation nette du
\emph{framerate} au-delà d'une vingtaine d'instances. Le nombre final a
été fixé à \textbf{50 arbres ``tree2''} et \textbf{140 palmiers}
(maillage bien plus léger), un compromis choisi délibérément pour
privilégier la \textbf{densité visuelle} de la forêt au prix d'un
\emph{framerate} réduit.

\section{Textures}

Plusieurs textures \texttt{cgp::opengl\_texture\_image\_structure} sont
utilisées :

\begin{itemize}
    \item \textbf{\texttt{texture\_rock}} et \textbf{\texttt{texture\_grass}}
    (avec une texture de sable intégrée comme texture principale du
    terrain) : le terrain mélange ces trois textures
    \emph{procéduralement} dans \texttt{mesh.frag.glsl}
    (\texttt{is\_terrain == 1}) en fonction de l'altitude
    \texttt{fragment.position.z} et de la pente
    \texttt{fragment.normal.z} :
    \begin{itemize}
        \item la roche (\texttt{texture\_rock}) est appliquée sur les
        pentes fortes (\texttt{slope} faible), via
        \texttt{smoothstep(0.150, 0.20, slope)} ;
        \item le sable apparaît près du niveau de la mer (basse
        altitude) ;
        \item l'herbe (\texttt{texture\_grass}) recouvre le reste.
    \end{itemize}
    Les trois contributions sont mélangées par interpolation
    (\emph{triplanar-like blending} basé sur la hauteur et la pente),
    chacune étant échantillonnée avec un facteur d'agrandissement
    différent (\texttt{uv*8}, \texttt{uv*6}, \texttt{uv*10}) pour éviter
    un effet de répétition trop visible ;
    \item \textbf{\texttt{texture\_sky}} : texture panoramique du ciel
    (\texttt{assets/sky.jpg}), utilisée à la fois pour la \emph{skybox}
    (section~\ref{sec:skybox}) et pour les reflets sur l'eau
    (section~\ref{sec:reflets}) ;
    \item les autres maillages (palmiers, arbres ``tree2'', bateau,
    épave, coffre, oiseaux) possèdent chacun leur(s) propre(s) texture(s)
    chargée(s) avec le maillage, appliquées via
    \texttt{image\_texture} et le modèle de Phong standard.
\end{itemize}

\section{Animation des oiseaux}
\label{sec:animation}

Une population d'oiseaux (\texttt{BirdGroup}) survole la scène. Chaque
groupe possède des paramètres d'orbite elliptique
(\texttt{orbit\_radius\_x}, \texttt{orbit\_radius\_y},
\texttt{base\_altitude}, \texttt{angular\_speed}, \texttt{phase}) : la
position de chaque oiseau à l'instant $t$ est calculée par
$x = r_x \cos(\omega t + \phi)$, $y = r_y \sin(\omega t + \phi)$,
$z = z_0$, ce qui produit une trajectoire elliptique fermée autour d'un
centre. Certains groupes (\texttt{forest\_bird = true}) effectuent en
plus un aller-retour entre un nid (\texttt{nest\_position}) situé dans
la forêt et leur orbite, sur une durée \texttt{tour\_duration}.

L'orientation de chaque oiseau (vecteur \texttt{bird\_direction}) est
déduite de la dérivée de sa trajectoire (direction du déplacement), et
un battement d'ailes est simulé en faisant osciller l'angle des ailes
(\texttt{bird\_wing\_angle}) selon une sinusoïde
(\texttt{bird\_wing\_amplitude}, \texttt{bird\_wing\_frequency}). Les
différentes parties de l'oiseau (corps, tête, bec, ailes, queue) sont
chacune des \texttt{mesh\_drawable} distincts, dessinés -- comme la
végétation -- via \texttt{draw\_part\_instanced} pour gérer
efficacement plusieurs dizaines d'oiseaux en quelques appels de dessin,
y compris dans la passe d'ombre (les oiseaux projettent une ombre sur
le sol/l'eau).

\section{Optimisations de performance}

Pour conserver un rendu interactif malgré la richesse de la scène
(île + crevasse, eau animée, ciel, ombres, $\sim$190 arbres instanciés,
oiseaux), plusieurs optimisations ont été appliquées :

\begin{itemize}
    \item réduction de la résolution des grilles de terrain et d'eau de
    $1000\times1000$ à $200\times200$ sommets (division par 25 du nombre
    de triangles) ;
    \item \emph{instancing} GPU pour les palmiers, les arbres ``tree2''
    et les oiseaux, ramenant des centaines d'appels de dessin à
    quelques appels par type d'objet ;
    \item mise à jour \emph{en place} des VBOs de l'eau
    (\texttt{vbo\_position.update}/\texttt{vbo\_normal.update}) au lieu
    d'une réallocation par image, supprimant une fuite mémoire GPU
    majeure ;
    \item réduction de la résolution de la \emph{shadow map} de
    $2048\times2048$ à $1024\times1024$.
\end{itemize}

Le nombre d'instances de végétation (140 palmiers, 50 arbres
``tree2'') a finalement été choisi en privilégiant la densité de la
forêt par rapport au \emph{framerate}, ce dernier restant néanmoins
acceptable pour une démonstration sur le GPU intégré utilisé pendant le
développement.

\section{Conclusion}

Ce projet met en \oe{}uvre l'ensemble de la chaîne de rendu vue en
cours : génération procédurale de géométrie par bruit de Perlin
(terrain, crevasse, vagues), modèle d'illumination de Phong, texturage
multi-couches, \emph{environment mapping}/\emph{skybox} pour le ciel et
ses reflets sur l'eau, \emph{shadow mapping} avec PCF, chargement et
séparation de maillages OBJ multi-objets, et \emph{GPU instancing} pour
peupler la scène d'une forêt et d'oiseaux animés. Les principales
difficultés -- fuite mémoire sur l'eau animée, séparation du modèle
d'arbre ``tree2'' livré en un seul fichier OBJ multi-objets, et coût en
sommets de ce modèle -- ont été résolues ou contournées par des choix
explicites de compromis qualité/performance, documentés ci-dessus.

\end{document}
